\documentclass[conference]{IEEEtran}

\usepackage{amsmath}
\usepackage{amssymb}
\usepackage{booktabs}
\usepackage{cite}
\usepackage{xspace}

\title{Hardware Architecture Section Preview}
\author{Anonymous Submission}

\begin{document}
\maketitle

\begin{abstract}
This standalone file previews only the hardware architecture section for the
MXFP4 Gated DeltaNet FPGA accelerator manuscript.
\end{abstract}

\section{Hardware Architecture}

\subsection{Overview}

The proposed architecture implements one Gated DeltaNet decode layer for
batch-one autoregressive inference. For each token step $t$, the kernel receives
token-dependent vectors
\[
q_t, k_t, v_t \in \mathbb{R}^{d}
\]
and scalar control values
\[
\beta_t, g_t \in \mathbb{R}.
\]
The kernel updates a recurrent state matrix
\[
S_t \in \mathbb{R}^{d \times d}
\]
and produces an output vector
\[
o_t \in \mathbb{R}^{d}.
\]
The design focuses on the decode primitive after the query, key, and value
projections have already been generated. The main architectural goal is to keep
$S_t$ resident in on-chip FPGA memory across token steps.

\subsection{Gated DeltaNet Recurrence}

The hardware follows the Gated DeltaNet update rule used by the software
reference model. First, the previous state is queried by the key vector:
\[
r_t = S_{t-1} k_t ,
\]
where $r_t$ is the predicted value retrieved from the old state. The prediction
error is then computed as
\[
e_t = r_t - v_t .
\]
The state is updated using a beta-scaled rank-one correction:
\[
S_t = S_{t-1} - \beta_t e_t k_t^{T}.
\]
Finally, the output vector is produced from the updated state:
\[
o_t = S_t q_t .
\]
Depending on matrix layout, the implementation may use the transposed form
$S_t^{T}q_t$ or $q_tS_t$, but the hardware operation is still a state-vector
product.

\subsection{Persistent State Storage}

The recurrent state matrix has fixed size $d \times d$ and does not grow with
sequence length. For a model with $H$ heads, the total stored state is
\[
H d^2
\]
elements. This fixed-size property allows the state to be stored in partitioned
on-chip memory banks.

The proposed design keeps $S_t$ persistent across token steps. Therefore, the
top-level interface only transfers the current token inputs and the output
vector. It does not transfer the full recurrent state every token. This changes
the main decode bottleneck from external state movement to on-chip state update
bandwidth.

\subsection{Five-Phase Token Dataflow}

Each token is processed using a fixed five-phase schedule:
\[
\text{prepare} \rightarrow \text{read} \rightarrow \text{update}
\rightarrow \text{writeback} \rightarrow \text{output}.
\]
The dominant operations are the state-vector product and the rank-one state
update. With key-direction parallelism $P_K$, a full pass over the $d \times d$
state costs approximately
\[
T_{\text{pass}} \approx \frac{d^2}{P_K}
\]
cycles. Since the design requires one main state read phase and one
writeback/update phase, the token latency can be approximated as
\[
T_{\text{token}} \approx \frac{2d^2}{P_K} + O(d).
\]
The lower-order $O(d)$ terms come from vector operations such as beta scaling,
residual computation, and output correction.

\subsection{Native MXFP4 Datapath}

The accelerator uses MXFP4 operands in the low-precision computation path. Each
block of $B$ elements shares one E8M0 scale, while each element is represented
as a four-bit E2M1 value. A block can be written as
\[
x_i \approx 2^{s_b}\hat{x}_i, \quad i \in b,
\]
where $s_b$ is the shared block exponent and $\hat{x}_i$ is the E2M1 element
value.

This representation reduces storage and arithmetic cost while preserving local
dynamic range. Unlike a dequantized design, the proposed datapath does not
convert MXFP4 operands into BF16 or FP32 before multiplication.

\subsection{LUT-Based E2M1 Multiplication}

For two E2M1 values $a$ and $b$, the multiplier computes
\[
p = a \cdot b
\]
directly using LUT logic. Each operand is decoded into sign, exponent, and
mantissa fields:
\[
a = (-1)^{s_a} 2^{e_a} m_a,
\]
\[
b = (-1)^{s_b} 2^{e_b} m_b.
\]
The product sign is
\[
s_p = s_a \oplus s_b,
\]
the exponent path performs
\[
e_p = e_a + e_b - \text{bias},
\]
and the mantissa path computes
\[
m_p = m_a \cdot m_b.
\]
The result is normalized and converted into a fixed-point partial product. This
keeps the inner multiply small enough to map efficiently into FPGA LUTs.

\subsection{Block-Exponent Alignment}

Because MXFP4 uses shared block scales, products from different blocks may have
different effective exponents. If two input blocks have exponents $s_a$ and
$s_b$, then the product block has effective exponent
\[
s_{\text{prod}} = s_a + s_b.
\]
Before accumulation, the partial product is shifted according to this exponent:
\[
\tilde{p}_i = p_i \ll \Delta s_i,
\]
where $\Delta s_i$ aligns the product to the accumulator scale. The aligned
partial products are accumulated into a fixed-point accumulator:
\[
A = \sum_i \tilde{p}_i.
\]
This stage preserves the block-scale information without widening the inner
multiply to BF16 or FP32.

\subsection{Parallelism Parameters}

The architecture is controlled by three main compile-time parameters:
\[
B,\quad P_K,\quad P_V.
\]
Here, $B$ is the MXFP4 block size, $P_K$ is the key-direction parallelism, and
$P_V$ is the value-direction parallelism. $B$ determines how many E2M1 values
share one block exponent. $P_K$ controls how many state columns are processed in
parallel. $P_V$ controls how many value or output lanes are updated together.

The default configuration is
\[
B = 32,\quad P_K = 16,\quad P_V = 8.
\]
These parameters determine the tradeoff between latency, resource utilization,
and scale-alignment overhead.

\subsection{Summary}

The architecture maps the Gated DeltaNet recurrence onto a persistent-state FPGA
datapath. The recurrent state $S_t \in \mathbb{R}^{d \times d}$ is stored on
chip, token updates are scheduled as fixed read-update-write phases, and MXFP4
arithmetic is implemented natively through LUT-based E2M1 multiplication and
block-exponent-aware accumulation. This design targets the two main costs of
decode: moving recurrent state and performing repeated state-vector products.

\end{document}
